\section{Limitations and Remaining Work}

The simulator's cost model is a set of formulas over synthetic demand, without
protocol, storage or queueing delays. The cluster pilots establish that
placement matters on real hardware; they do not establish that the simulator's
numbers are right. The two pilots also changed the transaction count and the
core count together, so the difference between them cannot be attributed to
either alone, and that count was chosen by convenience rather than by a power
calculation.

Three limitations sit on the learning side and they are the substantial ones.
No run had converged: held-out evaluation was still improving at the final
training step in every case, so the reported numbers are a floor rather than a
ceiling. Nothing can yet be claimed about the movement weight, because the
spread between two training seeds at the fastest setting is 4.70\,ms whilst
changing the weight from 0.30 to 0.10 at a fixed seed changes latency by
2.97\,ms, so the difference between seeds is larger than the effect being
studied. And between 79.9\% and 89.8\% of the actions the agent selects are
inapplicable in the state it selected them in, against 80.4\% for the random
baseline, so it is winning on the minority of its decisions that land.

That last figure is the largest measured inefficiency in the project and the
first thing to fix. It is a known failure mode as the illegal share of an
action list grows~\cite{huang2022masking}, and action-elimination networks
built on this kind of agent address it directly~\cite{zahavy2018elimination}.
After that, in order: more seeds per setting, which is cheap and unblocks the
movement-weight question; better credit assignment, since a move is charged in
full when it is made whilst the latency it buys arrives over later windows,
which $n$-step returns~\cite{hessel2018rainbow} would address; longer
training; separating the transaction count from the core count; building the
harness that would let the trained agent act on the real cluster, which does
not exist yet; and adding a cost-model rebalancer in the style of
Lion~\cite{lion2024} as a stronger opponent.
